\documentclass[12pt]{article}
\title{Asset Defender Risk Logic Q\&A Summary}
\author{}
\date{\today}
\begin{document}
\maketitle

\section*{Q\&A Summary}

1) Source of truth in production

- Live popup verdict: popup-computed finalRisk.
- Persisted/reporting history: server final_predictions.risk_score.
- Current state: not guaranteed identical.

2) Unify weighted formula (0.3/0.7 vs 0.6/0.4)

- Yes, unify everywhere.
- Use one shared formula constant across popup, server, and reporting UI.

3) Rename "Moderate" / "modern" wording

- Yes, rename for clarity.
- Recommended labels: Very Low, Low, Medium, High, Critical.

4) Hard-threshold docs in one table

- Yes.
- Include risk bands, category rules (hacker/scam/bot/suspicious/genuine), override rules, and boundary examples.

5) Owner of final classification logic

- Recommended production owner: server.
- Popup can show provisional output until server confirmation.
- If both paths remain, enforce parity tests.

6) Persist final category + exact reasons

- Yes, required for auditability.
- Persist category, score, triggered rule IDs, feature values, model/heuristic version, timestamp, and source.

7) Conflict policy when popup and server disagree

- Server wins for stored/reporting/enforcement decisions.
- UI shows provisional client result, then replaces with confirmed server result.
- Log mismatch diffs.

8) Regression tests for thresholds (39/40, 69/70, etc.)

- Yes.
- Add below/at/above boundary tests, override-path tests, insufficient-evidence tests, and parity tests.

\section*{Additional Questions Recommended}

1. Should popup scoring be production-disabled and used only for offline/debug mode?
2. Should every decision store a deterministic decision-trace payload for replay?
3. What server confirmation SLA should replace provisional client output?
4. Which threshold set is versioned as v1.0, and how do we roll back v1.1 safely?

\end{document}
